\subsection{Commutant and Duality}

Let A and B be k-algebras. Recall (III, §4, No.~3, p.~466) that a bimodule over the algebras A and B is an (A,B)-bimodule P on which the two k-module structures deduced from the A- and B-module structures coincide. To avoid any ambiguity, we say that P is an \((\mathrm{A},\mathrm{B})_{k}\) -bimodule. Let P be an \((\mathrm{A},\mathrm{B})_{k}\) -bimodule. We denote by \(P^{*}\) the dual \(\operatorname{Hom}_{\mathrm{A}}(\mathrm{P},\mathrm{A})\) of the left A-module underlying P. It is a \((\mathrm{B},\mathrm{A})_{k}\) -bimodule (II, §1, No.~14, p.~226); for \(a \in A\) , \(b \in B\) , \(x \in P\) , and \(x^{*} \in P^{*}\) , we have

\[
\langle x, b x ^ {*} a \rangle = \langle x b, x ^ {*} \rangle a.\tag{1}
\]

We denote by \(_{s}A_{d}\) the algebra A viewed as an \((\mathrm{A},\mathrm{A})_{k}\) -bimodule (loc. cit.) and by \(\Lambda:P\otimes_{B}P^{*}\to_{s}A_{d}\) the homomorphism of \((\mathrm{A},\mathrm{A})_{k}\) -bimodules determined by

\[
\Lambda (x \otimes x ^ {*}) = \langle x, x ^ {*} \rangle\tag{2}
\]

for \(x \in \mathrm{P}\) and \(x^{*} \in \mathrm{P}^{*}\) . We denote by \(\widetilde{\mathrm{P}}\) the dual \(\mathrm{Hom}_{\mathrm{B}}(\mathrm{P},\mathrm{B})\) of the right B-module underlying \(\mathrm{P}\) ; it is a \((\mathrm{B},\mathrm{A})_{k}\) -bimodule. We denote by \(\widetilde{\Lambda}: \widetilde{\mathrm{P}} \otimes_{\mathrm{A}} \mathrm{P} \to {}_{s}\mathrm{B}_{d}\) the homomorphism of \((\mathrm{B},\mathrm{B})_{k}\) -bimodules determined by

\[
\widetilde {\Lambda} (\widetilde {x} \otimes x) = \langle \widetilde {x}, x \rangle\tag{3}
\]

for \(x \in P\) and \(\widetilde{x} \in \widetilde{P}\) .

Now, suppose that the mapping \(b \mapsto b_{\mathrm{P}}\) is a bijection from B to \(\operatorname{End}_{\mathrm{A}}(\mathrm{P})\) ; it is then an isomorphism from B to the opposite algebra of \(\operatorname{End}_{\mathrm{A}}(\mathrm{P})\) . The canonical homomorphism of \(\mathbf{Z}\) -modules from \(\mathrm{P}^* \otimes_{\mathrm{A}} \mathrm{P}\) to \(\operatorname{End}_{\mathrm{A}}(\mathrm{P})\) (II, §4, No.~2, p.~271) then determines a homomorphism of \(\mathbf{Z}\) -modules \(\Theta: \mathrm{P}^* \otimes_{\mathrm{A}} \mathrm{P} \to \mathrm{B}\) defined by

\[
x \Theta (x ^ {*} \otimes y) = \langle x, x ^ {*} \rangle y\tag{4}
\]

for \(x, y \in \mathrm{P}\) and \(x^{*} \in \mathrm{P}^{*}\) . Because of (1), this homomorphism is (B,B)-linear, and we have

\[
\Theta (x ^ {*} \otimes y) y ^ {*} = x ^ {*} \langle y, y ^ {*} \rangle\tag{5}
\]

for \(y \in \mathrm{P}\) and \(x^{*}, y^{*} \in \mathrm{P}^{*}\) . From (4) and (5), we deduce the following equalities in the \((\mathrm{B}, \mathrm{B})_{k}\) -bimodule \(\mathrm{P}^{*} \otimes_{\mathrm{A}} \mathrm{P}\) :

\[
(y ^ {*} \otimes x) \Theta (x ^ {*} \otimes y) = y ^ {*} \otimes \langle x, x ^ {*} \rangle y = y ^ {*} \langle x, x ^ {*} \rangle \otimes y = \Theta (y ^ {*} \otimes x) (x ^ {*} \otimes y)
\]

for \(x, y \in \mathrm{P}\) and \(x^{*}, y^{*} \in \mathrm{P}^{*}\) , and therefore

\[
s \Theta (t) = \Theta (s) t\tag{6}
\]

for \(s, t \in \mathrm{P}^* \otimes_{\mathrm{A}} \mathrm{P}\) . Likewise, for \(x, y\) in \(\mathrm{P}\) and \(x^*, y^*\) in \(\mathrm{P}^*\) , we deduce the following equalities in the \((\mathrm{A}, \mathrm{A})_k\) -bimodule \(\mathrm{P} \otimes_{\mathrm{B}} \mathrm{P}^*\) from (4) and (5):

\[
(x \otimes x ^ {*}) \langle y, y ^ {*} \rangle = x \otimes \Theta (x ^ {*} \otimes y) y ^ {*} = x \Theta (x ^ {*} \otimes y) \otimes y ^ {*} = \langle x, x ^ {*} \rangle (y \otimes y ^ {*}),
\]

and therefore

\[
u \Lambda (v) = \Lambda (u) v\tag{7}
\]

for \(u, v \in \mathrm{P} \otimes_{\mathrm{B}} \mathrm{P}^{*}\) .

For any element \(x^{*}\) of \(P^{*}\) , denote the B-linear mapping \(x \mapsto \Theta(x^{*} \otimes x)\) from P to B by \(\sigma(x^{*})\) . We thus define a mapping \(\sigma\) from \(P^{*}\) to \(\widetilde{P}\) that is (B, A)-linear and satisfies, by definition,

\[
\Theta (x ^ {*} \otimes y) = \langle \sigma (x ^ {*}), y \rangle\tag{8}
\]

for \(x^{*} \in P^{*}\) and \(y \in P\) . By the definition of \(\widetilde{\Lambda}\) , we therefore have

\[
\Theta = \widetilde {\Lambda} \circ (\sigma \otimes 1 _ {\mathrm{P}}).\tag{9}
\]

Suppose that the mapping \(a \mapsto a_{\mathrm{P}}\) from A to \(\operatorname{End}_{\mathrm{B}}(\mathrm{P})\) is bijective; it is then an algebra isomorphism. Analogously, we define a homomorphism of \((\mathrm{A}, \mathrm{A})_k\) -bimodules \(\widetilde{\Theta}: \mathrm{P} \otimes_{\mathrm{B}} \widetilde{\mathrm{P}} \to {}_s \mathrm{A}_d\) by setting

\[
\widetilde {\Theta} (x \otimes \widetilde {y}) y = x \langle \widetilde {y}, y \rangle\tag{10}
\]

for \(x, y \in \mathrm{P}\) and \(\widetilde{y} \in \widetilde{\mathrm{P}}\) . We also define a homomorphism of \((\mathrm{B}, \mathrm{A})_k\) -bimodules \(\widetilde{\sigma}: \widetilde{\mathrm{P}} \to \mathrm{P}^*\) by setting

\[
\widetilde {\Theta} (x \otimes \widetilde {y}) = \langle x, \widetilde {\sigma} (\widetilde {y}) \rangle\tag{11}
\]

for \(x\in \mathrm{P},\widetilde{y}\in \widetilde{\mathrm{P}}\) . We have

\[
\widetilde {\Theta} = \Lambda \circ (1 _ {P} \otimes \widetilde {\sigma}).\tag{12}
\]

PROPOSITION 1. --- Suppose that the mappings \(b \mapsto b_{P}\) from B to \(\operatorname{End}_{\mathrm{A}}(\mathrm{P})\) and \(a \mapsto a_{P}\) from A to \(\operatorname{End}_{\mathrm{B}}(\mathrm{P})\) are bijective. Then \(\sigma\) and \(\widetilde{\sigma}\) are inverse isomorphisms, and we have the relations

\[
\Lambda = \widetilde {\Theta} \circ (1 _ {\mathrm{P}} \otimes \sigma),\tag{13}
\]

\[
\widetilde {\Lambda} = \Theta \circ (\widetilde {\sigma} \otimes 1 _ {\mathrm{P}}).\tag{14}
\]

For \(x \in \mathrm{P}\) , \(x^{*} \in \mathrm{P}^{*}\) , and \(y \in \mathrm{P}\) , by relations (4), (8), (10), and (11), we have

\[
\langle x, x ^ {*} \rangle y = x \Theta (x ^ {*} \otimes y) = x \langle \sigma (x ^ {*}), y \rangle = \widetilde {\Theta} (x \otimes \sigma (x ^ {*})) y = \langle x, \widetilde {\sigma} (\sigma (x ^ {*})) \rangle y.\tag{15}
\]

Likewise, for \(x \in P\) , \(\widetilde{y} \in \widetilde{P}\) , and \(y \in P\) , we have

\[
x \langle \widetilde {y}, y \rangle = \widetilde {\Theta} (x \otimes \widetilde {y}) y = \langle x, \widetilde {\sigma} (\widetilde {y}) \rangle y = x \Theta (\widetilde {\sigma} (\widetilde {y}) \otimes y) = x \langle \sigma (\widetilde {\sigma} (\widetilde {y})), y \rangle .\tag{16}
\]

Because of the assumptions, P is faithful as an A-module and as a B-module. From relations (15) and (16), respectively, we deduce \(\widetilde{\sigma} \circ \sigma = 1_{P^{*}}\) and \(\sigma \circ \widetilde{\sigma} = 1_{\widetilde{P}}\) . Relations (13) and (14) then follow from (12) and (9), respectively.

Remarks. --- 1) Suppose that the mapping \(b \mapsto b_{\mathrm{P}}\) from B to \(\operatorname{End}_{\mathrm{A}}(\mathrm{P})\) is bijective. Then

\begin{enumerate}
\def\labelenumi{\alph{enumi})}
\item
  the B-module P can be identified with the countermodule of P; it is therefore faithful and balanced;
\item
  the mapping \(a \mapsto a_{P}\) from A to \(\operatorname{End}_{\mathrm{B}}(\mathrm{P})\) is bijective if and only if the A-module P is faithful and balanced.
\end{enumerate}

\begin{enumerate}
\def\labelenumi{\arabic{enumi})}
\setcounter{enumi}{1}
\tightlist
\item
  Under the assumptions of Proposition 1, the A-module P is balanced; since the rings \(A_{P}\) and \(A_{P}^{\prime}\) have the same center (VIII, p.~77), there is an isomorphism \(\varphi\) from the center \(Z(A)\) of the ring A to the center \(Z(B)\) of the ring B, determined by the relation \(\varphi(z)_{\mathrm{P}} = z_{\mathrm{P}}\) for \(z \in \mathrm{Z}(A)\) . Moreover, the endomorphisms of the \((\mathrm{A}, \mathrm{B})_{k}\) -bimodule P are the homotheties \(z_{P}\) where z runs through \(Z(A)\) ; the automorphisms of the \((\mathrm{A}, \mathrm{B})_{k}\) -bimodule P are the homotheties \(z_{P}\) where z is invertible in \(Z(A)\) .
\end{enumerate}
% semantic-fragments: legacy-input-seam
\relax{}
